\question{What lies outside this report's technical scope?}
\label{sec:outside}
\statusOutside

\begin{shortanswer}
Four things, and they are not equivalent. Some are legal determinations that are not ours to make;
one depends on a consortium decision that has not been taken; one requires a deployment that does
not exist; and one belongs to a different work package. Each is set out below with the technical
facts counsel needs in order to decide it, so that the boundary is a handover rather than a gap.
\end{shortanswer}

\subsection{Legal determinations we do not make}

Whether a gradient \emph{is} personal data under Art.~4(1); whether the released model is anonymous;
the scope and status of the data protection impact assessment; and classification under the EU AI
Act. This report supplies inputs --- what is transmitted, who can see it, what attacks recover, what
the budget is --- and stops there.

We note only that the assessment's provisional recommendation, to treat gradients as personal and
potentially as health data until anonymisation is demonstrated, remains the technically correct
treatment of the \emph{default} configuration. What has changed is the set of configurations
available (\S\ref{sec:isolation}, \S\ref{sec:secagg}), not the status of the default.

\subsection{A decision the consortium has not yet taken}

The allocation of controller, joint-controller and processor roles under Art.~26/28 turns on
\textbf{who operates the SuperLink host} --- docport, or IKIM/KITE. That decision has not been made,
so the question is not merely legal but not yet decidable. The technical facts bearing on it:

\begin{itemize}
  \item Whoever operates the host is the party that, on the default configuration, receives
        individual attributable updates (\S\ref{sec:isolation}).
  \item With secure aggregation enabled, that party receives only the aggregate --- which may
        materially change the analysis, and is the strongest reason to take the configuration
        decision early rather than at Milestone~4.
  \item The node registry mapping public keys to admitted practices is centrally held by that
        operator, independently of the model path.
\end{itemize}

\subsection{Facts about a deployment that does not exist}

The assessment's fourth documented item --- administrative access, logging and export capabilities
on the SuperLink --- is the one item this report cannot close. We can document the full
\textbf{capability surface}, and Appendix~\ref{sec:config} does: the state database, log files and
rotation, run log retrieval, run listing and termination, and the node registry commands. What no
amount of engineering here can state is which of those an operator will enable, what retention they
will set, and who will hold administrative access. That is a fact about a host and an organisation,
not about this software.

Two further items sit here: the \textbf{real per-practice cohort sizes}, which are an MS3/MS5 fact
about the world and the input that decides whether a meaningful budget is reachable; and the
\textbf{practices' own IT environments}, which may introduce disclosure surfaces this project never
sees.

\subsection{A component this repository does not implement}

The federated analytics path. \S\ref{sec:aggdisclosure} answers its structural questions --- they
carry over unchanged, because the distinction is not architectural in Flower --- and states exactly
what specification would be needed to answer its disclosure-control questions. We flag it here as
well because silence on half of the assessment's framing would reasonably be read as a claim, and it
is not one.
